Building on advances in LLM-driven code generation, recent work has increasingly applied LLMs to the generation of high-performance kernels. To highlight the methodological patterns that have emerged across this landscape, the following sections review two dominant families of post-training techniques used to specialize LLMs for kernel generation: \emph{supervised fine-tuning} and \emph{reinforcement learning}.

\subsection{Supervised Fine-Tuning}

Supervised fine-tuning (SFT) has become a central methodology for enabling LLMs to synthesize high-quality kernels, relying on paired datasets that capture both high-level computational intent and low-level kernel implementation patterns. 
One influential line of work shows that the structure and clarity of model reasoning can strongly affect kernel correctness and performance. ConCuR~\cite{DBLP:journals/corr/abs-2510-07356} demonstrates this by constructing a curated dataset in which training samples are selected based on the conciseness of their reasoning processes, the speedup they achieve, and the diversity of their computational tasks. Fine-tuning on such data leads to KernelCoder, a model capable of generating CUDA kernels with state-of-the-art reliability and efficiency. Another direction builds paired training corpora through compiler alignment, where kernel implementations are automatically generated to mirror high-level operators. KernelLLM~\cite{kernelllm2025} adopts this strategy by using the Triton compiler to produce aligned PyTorch–Triton examples and by applying instruction tuning with structured prompts that explicitly encode the mapping between computation and kernel structure.  Together, these approaches show that well-designed supervised datasets can effectively specialize LLMs for robust and high-performance GPU kernel synthesis.


\subsection{Reinforcement Learning}

Reinforcement learning enhances kernel generation via iterative feedback. Kevin~\cite{baronio2025kevin} models kernel generation as a multi-turn optimization using cross-turn reward attribution for long-horizon credit assignment. QiMeng-Kernel~\cite{zhu2025qimengkernel} further structures optimization by applying RL hierarchically to macro-thinking strategies rather than low-level implementation. Recent approaches focus on robust reward mechanisms and verifiable evaluation. AutoTriton~\cite{li2025autotriton} addresses reward sparsity by combining structural assessments of generated kernels with execution-based runtime rewards, while TritonRL~\cite{woo2025tritonrl} extends this line of work through hierarchical reward decomposition and explicit verification of code outputs and intermediate reasoning traces. CUDA-L1 introduces contrastive RL with an LLM-as-a-judge for dense feedback, and is refined by CUDA-L2~\cite{su2025cudal2surpassingcublasperformance} to surpass cuBLAS performance. Finally, AscendKernelGen~\cite{cao2026ascendkernelgen} expands the preference learning paradigm to Ascend NPUs, combining CoT-based SFT with preference learning.


% Reinforcement learning (RL) has emerged as a powerful paradigm for improving kernel generation by enabling iterative optimization driven by execution-based and verification-based feedback. SparseRL~\cite{anonymous2026sparserl} exemplifies this approach by using compilation outcomes and execution characteristics as reward signals, leading to substantial improvements in sparse matrix–vector and sparse matrix–matrix multiplication tasks.  Kevin~\cite{baronio2025kevin} further formulates kernel synthesis as a multi-turn optimization process, in which the model repeatedly proposes candidate kernels and refines subsequent generations according to compilation feedback. To address the challenge of long-horizon credit assignment, this framework introduces cross-turn reward attribution and context compression mechanisms, enabling effective optimization over extended interaction trajectories.  Furthermore, QiMeng-Kernel~\cite{zhu2025qimengkernel} proposes a hierarchical framework including macro thinking and micro coding, decoupling optimization strategy from low-level implementation details, where the macro thinking is trained by reinforment learning to maximize hardware utilization.

% Recent advances also emphasize the importance of robust reward design and verifiable evaluation, particularly in settings where reward signals are sparse. AutoTriton~\cite{li2025autotriton} addresses this issue by combining structural assessments of generated kernels with execution-based runtime rewards, encouraging improvements in both algorithmic validity and empirical efficiency. TritonRL~\cite{woo2025tritonrl} extends this line of work through hierarchical reward decomposition and explicit verification of code outputs and intermediate reasoning traces, reducing susceptibility to reward hacking and promoting more faithful optimization behavior. CUDA-L1~\cite{li2025cuda} employs contrastive reinforcement learning to compare multiple candidate kernels and infer relative performance differences, while introducing an LLM-as-a-judge component that critiques intermediate outputs and provides dense, informative feedback when compilation or execution fails. CUDA-L2~\cite{su2025cudal2surpassingcublasperformance} builds upon CUDA-L1 by introducing more sophisticated reinforcement learning techniques and demonstrates superior performance compared to the widely-used industrial cuBLAS library. And AscendKernelGen~\cite{cao2026ascendkernelgen} further investigates kernel generation for Ascend NPUs by synthesizing a large-scale dataset incorporating chain-of-thought reasoning for SFT, followed by preference learning to further reinforce model capabilities.

